\documentclass[12pt]{article}
\usepackage[utf8]{inputenc}
\usepackage[T1]{fontenc}
\usepackage{hyperref}
\title{Elotepec Zapotec Language Resources for K-12 Curriculum Development}
\author{Compiled by Claude Sonnet 5.5 (claude-sonnet-5-5) for Andr\'as Kornai}
\date{}
\begin{document}
\maketitle

\section{Introduction}
Elotepec Zapotec (ISO 639-3 zte; also called Papabuco or Zapoteco de San Juan Elotepec) is spoken in San Juan Elotepec, Oaxaca, Mexico. Wikipedia gives about 200 speakers (1990 census). Wikitongues rates it EGIDS 7 (shifting) and lists it as unwritten. No school materials were found.

\section{Primary Educational Resources}
None found in search results.

\section{Linguistic and Reference Materials}
\begin{itemize}
\item Wikipedia article on Elotepec Zapotec. It mentions an online archive of the language funded by INALI and the NSF Documenting Endangered Languages program. The archive itself was not located. \url{https://en.wikipedia.org/wiki/Elotepec_Zapotec}
\item Wikitongues page for zte. \url{https://wikitongues.org/languages/zte/}
\end{itemize}

\section{Religious and Community Resources}
\begin{itemize}
\item Joshua Project page on the Elotepec Zapoteco people. \url{https://joshuaproject.net/people_groups/14699}
\end{itemize}

\section{Digital and Audio Resources}
\begin{itemize}
\item Global Recordings Network audio Bible stories in Elotepec Zapotec. \url{https://globalrecordings.net/en/language/zte}
\end{itemize}

\section{Curriculum Development Context}
The language has few speakers, who are mostly older adults, and it has no established writing system. Documentation and orthography development would come before any curriculum. The archive mentioned on Wikipedia should be tracked down as a possible source of recordings.

\end{document}